Haldane predicted a gap for integer-spin antiferromagnetic Heisenberg
chains, in contrast to their half-integer counterparts~\cite{Haldane1983Continuum,Haldane1983Nonlinear}.
For spin one, numerical work places the bulk gap near
$0.4105$~\cite{White1993Numerical}. Unlike the
Affleck--Kennedy--Lieb--Tasaki model~\cite{Affleck1987Rigorous}, the
Heisenberg chain is not frustration free: neighboring bonds cannot
simultaneously attain their singlet ground states. Finite-size criteria
based on frustration freeness~\cite{Knabe1988Energy,Gosset2016Local}
therefore do not apply directly.

OpenAI proved a unique ground state and a uniform gap
$\Delta_L>(\log20)/784$ for every even periodic length $L\ge2304$
using certified finite partition-function bounds~\cite{OpenAI2026Periodic}.
Here we explain the proof using tensor-network diagrams. We extend
it to other spin models under a condition that we identify as
reflection positivity, and give Lean proofs. We also improve the
spin-1 gap bound using OpenAI's finite trace calculations and trial
states; see details~\cite{SM}.

\emph{Thermal purity and the gap.---}
Consider
\begin{equation}
 H_L=\sum_{j=1}^{L}\bm S_j\cdot\bm S_{j+1},
 \qquad Z_L(\beta)=\Tr e^{-\beta H_L},
 \label{eq:core-model}
\end{equation}
with periodic indices, spin one, even $L\ge4$, and exchange coupling
one. Order energies with multiplicity and put
$\Delta_L=E_1(L)-E_0(L)$, illustrated by \levels.
For $\beta>0$, the Gibbs state $\rho_L=e^{-\beta H_L}/Z_L(\beta)$
has ordered weights $w_0\ge w_1\ge\cdots$. Its thermal purity
$\Pt(L,\beta)=\Tr\rho_L^2\le w_0$ gives
\begin{equation}
 e^{-\beta\Delta_L}=\frac{w_1}{w_0}
 \le\frac{1-w_0}{w_0}
 \le\frac{1-\Pt(L,\beta)}{\Pt(L,\beta)}.
 \label{eq:core-purity-gap}
\end{equation}
Purity close to one forces a single dominant Gibbs weight,
\specbars{.97,.02,.01}. Thus $\Pt>1/2$ forces a unique ground state, and
$1-\Pt\le q<1/2$, with $q>0$, gives
$\Delta_L\ge\beta^{-1}\log[(1-q)/q]$.

\emph{Proof strategy.---}
To keep the gap bound positive as we lower the temperature,
$\log(1/q)$ must grow proportionally to $\beta$.
Squaring and normalizing the Gibbs state doubles $\beta$, but does
not extend the chain. The key idea from the OpenAI proof is to construct a spatial state whose square doubles
the length, so that we can iterate both operations:
\begin{proofdiagram}
 \strategystairs
\end{proofdiagram}
% The diagram shows the intended small-error behavior, with constant factors suppressed.
If each step squares the purity errors, their
negative logarithms grow as $2^j$, matching $\beta_j=2^j\beta_0$.
Taking the logarithm and dividing by $\beta_j$ therefore leaves a
positive constant in the gap bound.
We must establish this estimate and cover the lengths between steps.
First we construct the spatial state.

\emph{Spatial transfer matrix.---}
Set $\tau=\beta/M$ and let $S_L(\tau)$ be the product of
$I-\tau h_{j,j+1}$ over even bonds followed by odd bonds, where
$h=\bm S\cdot\bm S$. Contracting the resulting finite network in
time or space gives~\cite{Suzuki1985Transfermatrix}
\begin{equation}
 \underbrace{\Tr S_L(\beta/M)^M}_{\bricknetwork{1}}
 =Z_L^{(M)}(\beta)
 =\underbrace{\Tr T_{\beta,M}^{L/2}}_{\bricknetwork{2}}.
 \label{eq:core-two-readings}
\end{equation}

In the Cartesian basis $\ket{x},\ket{y},\ket{z}$, the generators
$G_\alpha=iS^\alpha$ are real antisymmetric matrices. Hence
\begin{equation}
 -h=\sum_{\alpha=x,y,z}G_\alpha\otimes G_\alpha,
 \qquad
 I-\tau h=\sum_{\gamma=0,x,y,z}A_\gamma\otimes A_\gamma,
 \label{eq:core-bond}
\end{equation}
where $A_0=I$ and $A_\alpha=\sqrt{\tau}\,G_\alpha$.
Diagrammatically, $\splitgate=\sum_\gamma\bondends$.
Matching real bond factors express reflection
positivity~\cite{Dyson1978Phase,Frohlich1978Phase,Jaffe2016Characterization}.

For bond-label sequences $a,b\in\{0,x,y,z\}^M$, define the real matrix
$\mathcal F_{ba}=\Tr[(A_{b_1}A_{a_1})\cdots(A_{b_M}A_{a_M})]$.
The two sites sharing an even bond have identical factor orderings.
Contracting their shared sequence gives
\begin{equation}
 T_{\beta,M}=\mathcal F^{\mathsf T}\mathcal F\ge0,
 \qquad
 \tnchain{\mathcal F^{\mathsf T},\mathcal F}
 =\tnchain{T_{\beta,M}}.
 \label{eq:core-gram}
\end{equation}
This transfer matrix advances by two sites and is independent of $L$.

For even $n\ge4$, normalize $T_{\beta,M}^{n/2}$ by its trace to obtain
the density operator $\sigma_n^{(M)}$. The trace is nonzero because
$\mathcal F_{00}=3$ for the all-zero sequences.
Nonnegative transfer eigenvalues give spatial purity, squaring, and
power-sum monotonicity inequalities.
At each fixed $L,\beta$, the matrix product limit
$S_L(\beta/M)^M\to e^{-\beta H_L}$ passes these closed scalar
inequalities to exact partition functions.
We take this limit before the doubling argument, requiring neither a limiting
transfer matrix nor a Trotter error uniform in $L$.
The finite-Trotter section of the Supplemental Material~\cite{SM} proves these
statements. Thermal identities below concern the exact Gibbs operator,
not finite Trotter products.

\emph{Coupling and squaring.---}
The two purities are
\begin{align}
 \Ps(n,\beta)
 &=\lim_{M\to\infty}\Tr[(\sigma_n^{(M)})^2]
 =\frac{Z_{2n}(\beta)}{Z_n(\beta)^2},\nonumber\\
 \Pt(n,\beta)
 &=\Tr\rho_n(\beta)^2
 =\frac{Z_n(2\beta)}{Z_n(\beta)^2}.
 \label{eq:core-purities}
\end{align}
Both lie in $(0,1]$. Their contractions,
$\tnchain[ring]{\sigma_n^{(M)},\sigma_n^{(M)}}$ and
$\tncolumn[ring]{\rho_n,\rho_n}$, join normalized copies in
space and time. The rectangle $\Zt21$ denotes $Z_n(\beta)$, with width
$n$ and height $\beta$. Joining four copies in either order gives
\begin{proofdiagram}
 \torusgluing
\end{proofdiagram}
The two paths therefore give the exact identity
\begin{equation}
 \Ps(n,2\beta)\Pt(n,\beta)^2
 =\Pt(2n,\beta)\Ps(n,\beta)^2.
 \label{eq:core-coupling-s}
\end{equation}
It holds without reflection positivity; positivity is needed to use
its factors as purities.

For probabilities $u_i$, write $P=\sum_i u_i^2$. Their normalized
squares have purity $P'$, with
\begin{equation}
 1-P'=\frac{2\sum_{i<j}(u_i u_j)^2}{P^2}
 \le\frac{(1-P)^2}{2P^2}=f(1-P),
 \label{eq:core-square}
\end{equation}
where $f(z)=z^2/[2(1-z)^2]$.
Thus squaring concentrates weights,
$\specbars{.5,.3,.2}\mapsto\specbars{.66,.24,.1}$.
The same operation acts in the two directions:
\begin{align}
 \tnchain{\sigma_n^{(M)},\sigma_n^{(M)}}
 &=\Tr[(\sigma_n^{(M)})^2]\,\tnchain{\sigma_{2n}^{(M)}},\notag\\
 \tncolumn{\rho_n(\beta),\rho_n(\beta)}
 &=\Pt(n,\beta)\,\tncolumn{\rho_n(2\beta)}.
 \label{eq:core-operator-square}
\end{align}

Suppose $1-\Ps(n,\beta)\le p$ and $1-\Pt(2n,\beta)\le q$.
Dropping the favorable denominator $\Pt(n,\beta)^2\le1$ in
Eq.~\eqref{eq:core-coupling-s}, then squaring in space, gives $p'$.
Repeating the identity at length $2n$, then squaring in time, gives
$q'$:
\begin{align}
 p'&=f(1-(1-p)^2(1-q)),\nonumber\\
 q'&=f(1-(1-q)^2(1-p')).
 \label{eq:core-update}
\end{align}
For $0\le p,q,p',q'<1$, these bound
$1-\Ps(2n,2\beta)$ and $1-\Pt(4n,2\beta)$.
The thermal update uses the new spatial bound:
\begin{proofdiagram}
 \pbcsteps
\end{proofdiagram}

\emph{The finite seed and uniform gap.---}
The existing computation supplies an independent input.

\begin{lemma}[Finite seed~\cite{OpenAI2026Periodic}]
 \label{lem:arxiv-seed36}
 At $\beta_0=49/4$, the periodic spin-1 chain satisfies
 $\Ps(36,\beta_0)\ge1-47916/10^6$ and
 $\Pt(72,\beta_0)\ge1-181521/10^6$.
\end{lemma}

Accepting this lemma suffices for the following argument.
Six rationally verified, upward-rounded updates give
$p_6=0.0084709$ and $q_6=0.0055149$ at
$(n_6,\beta_6)=(2304,784)$, as shown in
Fig.~\ref{fig:core-region}. The doubling-argument section of the Supplemental Material~\cite{SM} records the exact orbit and rounding rule.

For $p,q\le\varepsilon\le0.01$, both updates are at most
$5\varepsilon^2$; the Supplemental Material~\cite{SM} proves this scalar estimate.
Restarting at $(2304,784)$ gives bounds
$\varepsilon_j=20^{-2^j}/5$ at
$n_j=2304\,2^j$, $\beta_j=784\,2^j$.
To cover every even $L\ge2304$, choose $j$ with
$n_j\le L\le2n_j$. Spatial power-sum monotonicity at $\beta_j$
and $2\beta_j$ yields
\begin{align}
 \Pt(L,\beta_j)
 &\ge[\Pt(2n_j,\beta_j)\Ps(n_j,\beta_j)^2]^{L/(2n_j)}
 \nonumber\\
 &\ge(1-\varepsilon_j)^3\ge1-3\varepsilon_j.
 \label{eq:core-interpolation}
\end{align}
Consequently the ground state is unique. Moreover,
$3\varepsilon_j/(1-3\varepsilon_j)
\le(60/97)20^{-2^j}<20^{-2^j}$.
Equation~\eqref{eq:core-purity-gap} therefore proves
\begin{equation}
 \Delta_L>\frac{\log20}{784}\simeq0.003821
 \qquad\text{for every even }L\ge2304.
 \label{eq:core-gap}
\end{equation}

\emph{Finite seed bounds.---}

The starting lemma requires two ratios at $\beta_0=49/4$:
\begin{align}
 \textcolor{mpored}{\Ps(36,\beta_0)}
 &=\frac{Z_{72}(\beta_0)}{Z_{36}(\beta_0)^2},\notag\\
 \textcolor{mpoblue}{\Pt(72,\beta_0)}
 &=\frac{Z_{72}(2\beta_0)}{Z_{72}(\beta_0)^2}.
 \label{seed:ratios}
\end{align}
\begin{figure}[t]
 \centering
 \begin{tikzpicture}
  \begin{axis}[
   width=\columnwidth,height=5.3cm,
   xmin=0,xmax=.23,ymin=0,ymax=.24,
   axis lines=left,
   axis line style={black!60,line width=.4pt},
   tick align=outside,
   xlabel={$p\ge1-\Ps(n,\beta)$},
   ylabel={$q\ge1-\Pt(2n,\beta)$},
   xlabel style={text=mpored},
   ylabel style={text=mpoblue},
   xtick={0,.1,.2},ytick={0,.1,.2},
   xticklabels={$0$,$0.1$,$0.2$},
   yticklabels={$0$,$0.1$,$0.2$},
   clip=false]
   \addplot[fill=teal!12,draw=teal!70!black,line width=.5pt]
    coordinates {
     (0,0) (0,.2006) (.01,.1975) (.02,.194) (.03,.1902)
     (.04,.186) (.05,.1813) (.06,.1762) (.07,.1706)
     (.08,.1645) (.09,.1579) (.10,.1506) (.11,.1427)
     (.12,.1341) (.13,.1247) (.14,.1145) (.15,.1035)
     (.16,.0915) (.17,.0785) (.18,.0644) (.19,.0491)
     (.20,.0326) (.21,.0146) (.2176,0) (0,0)};
   \fill[black!35] (axis cs:0,0) rectangle (axis cs:.01,.01);
   \addplot[black,mark=*,mark size=1.3pt,line width=.7pt]
    coordinates {
     (.0479160,.1815210)
     (.0605003,.1733838)
     (.0686630,.1632461)
     (.0713693,.1447318)
     (.0633146,.1055663)
     (.0376144,.0446518)
     (.0084709,.0055149)};
   \node[anchor=east,inner sep=2pt,align=right]
    at (axis cs:.0479160,.1815210) {seed\\$36$};
   \node[anchor=west,inner sep=2pt]
    at (axis cs:.0713693,.1447318) {$288$};
   \node[anchor=west,inner sep=2pt]
    at (axis cs:.0084709,.0055149) {$2304$};
   \node[text=teal!70!black,align=center]
    at (axis cs:.142,.071) {numerical\\convergence region};
  \end{axis}
 \end{tikzpicture}
 \caption{The asymmetric recurrence~\eqref{eq:core-update} and its
 spin-1 seed orbit at $n_k=36\,2^k$ and $\beta_k=(49/4)2^k$. Six
 rationally verified, upward-rounded steps carry the seed at $n_0=36$
 into the analytical target square $[0,0.01]^2$ (gray) at $n_6=2304$;
 the teal region shows the numerical basin of attraction.}
 \label{fig:core-region}
\end{figure}

Each numerator needs a lower bound; each denominator needs an upper
bound. Exact trace calculations on $4$--$12$ sites and a $72$-site
trial state supply them~\cite{OpenAI2026Periodic}:
\begin{center}
\setlength{\tabcolsep}{3pt}
\begin{tabular}{@{}clll@{}}
\toprule
 & Quantity & Bound & Method\\
\midrule
$\Zt{4}{1}$ & $Z_{72}(\beta_0)$ & lower & MPS + small-ring traces\\
$\Zt{2}{1}$ & $Z_{36}(\beta_0)$ & upper & MPS + small-ring traces\\
$\Zt{4}{2}$ & $Z_{72}(2\beta_0)$ & lower & MPS\\
$\Zt{4}{1}$ & $Z_{72}(\beta_0)$ & upper & small-ring traces\\
\bottomrule
\end{tabular}
\end{center}
Readers may skip this explanation and use Lemma~\ref{lem:arxiv-seed36} directly.
Defining decimals below are exact rationals; approximate values are
only explanatory.

\paragraph{Step 1: Spatial trace moments.}
A continuous-time, one-site spatial transfer operator gives
\begin{equation}
 Z_n(\beta)=\Tr K_\beta^n
 =\tnchain[node=xvertex,ring]{,\cdots,}_{\,n},
 \qquad n\ge4.
 \label{seed:transfer}
\end{equation}
It is self-adjoint but need not be positive, unlike the positive
two-site transfer matrix used in the doubling argument.
Anticipating the trial energy below, set $\kappa=700741/500000$ and
\begin{align}
 \widehat Z_n(\beta)&=e^{-n\beta\kappa}Z_n(\beta),\notag\\
 \widehat K&=e^{-\beta_0\kappa}K_{\beta_0},\qquad
 \widehat Z_n(\beta_0)=\sum_i\lambda_i^n.
 \label{seed:shift}
\end{align}
The $\lambda_i$ are real eigenvalues of $\widehat K$, possibly negative.
The shift cancels from both purities. Changing $n$ changes only the
power: short-chain traces constrain the same spectrum needed at
$36$ and $72$ sites.

\paragraph{Step 2: Separate the moments by symmetry.}
Spin rotations preserve the Heisenberg Hamiltonian and induce
transfer symmetries. The commuting $\pi$ rotations split four sign
sectors: an invariant sector $J$ and three sectors with identical
spectra, denoted $N$, exchanged by axis permutations. Thus
\begin{align}
 \seedring{0}\quad \widehat Z_n&=J_n+3N_n,\notag\\
 \seedring{1}\quad \widehat Z_n(P)&=J_n-N_n.
 \label{seed:sectors}
\end{align}
Here $J_n=\sum_{i\in J}\lambda_i^n$, and $N_n$ sums one $N$ list;
all traces are at $\beta_0$.
The twist $P$ rotates one end of the physical closing bond by $\pi$,
rather than inserting a global rotation into the untwisted physical
trace.

Integer trace recursions with explicit truncation and rounding bounds
enclose both traces for $4\le n\le12$.
Inverting Eq.~\eqref{seed:sectors} gives sector moments, including
\begin{equation}
 J_{12}\le m_J=1.0657205,\qquad
 N_{12}\le m_N=0.2783155.
 \label{seed:masses}
\end{equation}

\paragraph{Step 3: A trial state gives a lower bound.}
Integer tensors of bond dimension $26$ define a periodic $72$-site
MPS $\psi$. Its norm and one-bond expectation are exact contractions:
\begin{align}
 \mathcal N&=\langle\psi,\psi\rangle=\seedmps{0},\notag\\
 \mathcal E&=\langle\psi,h_{1,2}\psi\rangle=\seedmps{1}.
 \label{seed:trial}
\end{align}
Each diagram contains $72$ tensors per layer. Translation invariance gives
$\langle\psi,H_{72}\psi\rangle/(72\mathcal N)
=\mathcal E/\mathcal N<-\kappa$.
The variational principle gives $E_0(72)<-72\kappa$.
The shift adds $\kappa$ on each of the $72$ sites, so the shifted
ground energy is negative and its Boltzmann weight exceeds one:
\begin{equation}
 \widehat Z_{72}(\beta)
 \ge e^{-\beta(E_0(72)+72\kappa)}>1
 \qquad(\beta>0).
 \label{seed:trial-lower}
\end{equation}
This calculation supplies both
$\sum_i\lambda_i^{72}>1$ at $\beta_0$ and the thermal numerator
$\widehat Z_{72}(2\beta_0)>1$.

\paragraph{Step 4: A lower bound on the leading spatial eigenvalue.}
Suppose every $|\lambda_i|\le\ell=0.999$.
Write $u_i=\lambda_i^{12}\ge0$, so $\lambda_i^{72}=u_i^6$.
In $J$, the total mass is at most $m_J$ and each entry is at most
$\ell^{12}$. Convexity increases the sixth-power sum when mass
passes from a smaller entry to a larger one, up to the cap.
Since $\ell^{12}<m_J<2\ell^{12}$, the extremum has one entry
$\ell^{12}$ and remainder $m_J-\ell^{12}$:
\begin{proofdiagram}
 \seedweightbars
\end{proofdiagram}
Taking sixth powers suppresses the remainder.
Each $N$ list contributes at most $m_N^6$, giving
\begin{align}
 \widehat Z_{72}(\beta_0)
 &\le\ell^{72}+(m_J-\ell^{12})^6+3m_N^6\notag\\
 &<0.932<1.
 \label{seed:dominance}
\end{align}
This contradicts Eq.~\eqref{seed:trial-lower}.
Hence some $|\lambda_0|>\ell$; it belongs to $J$ because
$m_N<\ell^{12}$.

\paragraph{Step 5: Polynomial bounds on the remaining spectrum.}
A nonnegative polynomial $F(x)=\sum_{j=4}^{12}f_jx^j$ converts known
sector moments into individual eigenvalue bounds:
\begin{equation}
 F(\lambda_i)\le\sum_{a\in\mathcal S}F(\lambda_a)
 =\sum_{j=4}^{12}f_j\mathcal S_j\le B,
 \qquad i\in\mathcal S.
 \label{seed:filter}
\end{equation}
Here $\mathcal S$ is $J$ or $N$.
Where $F(x)>B$, no eigenvalue can occur. Two filters give
$|\lambda|<U=1.00139$ in $J$ and $|\lambda|<V=0.872$ in $N$.
The Supplemental Material~\cite{SM} gives their coefficients and checks both
exterior rays.

Removing $\lambda_0$ leaves $J$ mass at most $m_J-\ell^{12}$.
Since $V^{12}<m_N<2V^{12}$, convexity bounds each $N$ list by
one entry at $V^{12}$ and the remainder.
For $k=36,72$, the residual moment $\sum_{i\ne0}\lambda_i^k$
is therefore at most
\begin{align}
 R_k={}&(m_J-\ell^{12})^{k/12}\notag\\
 &+3\left[V^k+(m_N-V^{12})^{k/12}\right].
 \label{seed:residual}
\end{align}
Retain $\lambda_0^{72}$ in the spatial numerator and bound its
denominator by $(\lambda_0^{36}+R_{36})^2$.
For thermal purity, use the trial numerator and
$\widehat Z_{72}(\beta_0)\le U^{72}+R_{72}$. Thus
\begin{align}
 \Ps(36,\beta_0)&\ge[1+R_{36}/\ell^{36}]^{-2}=1-p,\notag\\
 \Pt(72,\beta_0)&\ge[U^{72}+R_{72}]^{-2}=1-q.
 \label{seed:long}
\end{align}
Exact rational comparisons give the starting bounds
\begin{equation}
 p\le\frac{47916}{10^6},\qquad q\le\frac{181521}{10^6}.
 \label{seed:bounds}
\end{equation}
For orientation, $p\simeq0.0479157710062$ and
$q\simeq0.181520168033$; these approximations are not recurrence inputs.

The same trace moments, the existing $120$-site trial state, and
temperature lifting give $\Delta_L>0.07290$ for every even $L\ge36$.
The lifted bound covers $L\ge60$; an unlifted estimate bridges
$36\le L\le72$. The Supplemental Material~\cite{SM} gives the proof, and the
\href{https://github.com/LionSR/CertifiedGaps}{repository} provides
the certificate.


\emph{Reflection-positive models.---}
The argument requires locality and reflection
positivity~\cite{Frohlich1978Phase,Jaffe2016Characterization}, not
spin-rotation symmetry. Consider a translation-invariant ring of $L$
cells with Hamiltonian
\begin{equation}
 H_L=\sum_{j=1}^L h_{j,j+1}+\sum_{j=1}^L o_j,
 \qquad j+L\equiv j.
 \label{scope:hamiltonian}
\end{equation}
Here $h$ couples neighboring cells and $o$ acts within a cell.
Let the real matrix $F=F^{\mathsf T}$, $F^2=I$, represent reflection
within a cell; for a dimer, $F$ exchanges the two spins.
We take $o$ real symmetric and require $FoF=o$, so the cell
Hamiltonian respects this reflection. Across a bond, we require
\begin{equation}
 -h=\sum_\alpha C_\alpha\otimes FC_\alpha F,
 \label{scope:rp}
\end{equation}
where each $C_\alpha$ is real and either symmetric or antisymmetric.
This pairs each bond factor with its reflected copy. Contracting two
neighboring cells therefore gives a positive Gram transfer matrix,
as in Eq.~\eqref{eq:core-gram}; see details~\cite{SM}.
We can then use the same spatial purity inequalities and doubling
argument as for the spin-1 chain.

An admissible seed at even $n_0$ then gives a unique ground state
and a uniform gap for every even $L\ge n_0$, using finite-prefix
interpolation and an analytic tail bound%
~\leanmark{https://sirui-lu.com/CertifiedGaps/docs/CertifiedGaps/Ring/Gap.html\#CertifiedGaps.RPRing.gapGT_of_chainCertificate}.
Table~\ref{tab:models} includes staggered-field chains~\cite{Oshikawa1997Field},
bond-alternating chains~\cite{Barnes1999Alternating,Yamamoto1995Elementary},
chains with single-ion anisotropy~\cite{Chen2003Ground}, and Heisenberg
ladders~\cite{Dagotto1996Surprises}. Each model needs its own finite
seed; the analytical argument is unchanged.
\begin{table*}[t]
\caption{Uniform gap bounds for reflection-positive models. Every even
length at or above the stated threshold is covered. A cell or rung
contains two spins; $L$ counts the units shown. Spin operators have
$\bm S^2=S(S+1)I$; $\sigma^\alpha$ are Pauli matrices. Numerical values
are finite-ring estimates on $12$--$20$ spins, not certified uniform
bounds, except for the spin-1 Heisenberg bulk estimate~\cite{White1993Numerical}.
The mixed-field Ising row gives the kernel-checked bound.}
\label{tab:models}
\setlength{\tabcolsep}{4pt}
\renewcommand{\arraystretch}{1.25}
\begin{tabular}{@{}p{.19\textwidth}p{.43\textwidth}lll@{}}
\toprule
Model & Hamiltonian & Even lengths & Gap $\ge$ & Numerical\\
\midrule
Staggered field, $S=\tfrac12$
 & $\sum_j\bm S_j\cdot\bm S_{j+1}-\sum_j(-1)^jS_j^z$
 & $L\ge8$ sites & $1.1489$ & $1.874$\\
Dimerized, $S=\tfrac12$
 & $\sum_j[1+(-1)^j/2]\,\bm S_j\cdot\bm S_{j+1}$
 & $L\ge4$ cells & $0.5347$ & $1.189$\\
Dimerized, $S=1$
 & $\sum_j[1+9(-1)^j/10]\,\bm S_j\cdot\bm S_{j+1}$
 & $L\ge2$ cells & $0.5288$ & $1.763$\\
Single-ion anisotropy, $S=1$
 & $\sum_j\bm S_j\cdot\bm S_{j+1}+4\sum_j(S_j^z)^2$
 & $L\ge12$ sites & $0.8994$ & $2.269$\\
Heisenberg ladder, \mbox{$S=\tfrac12$}
 & $\sum_{j,a=1,2}\bm S_{j,a}\cdot\bm S_{j+1,a}
    +4\sum_j\bm S_{j,1}\cdot\bm S_{j,2}$
 & $L\ge4$ rungs & $1.1024$ & $3.138$\\
Ising ladder
 & $-\sum_j(\sigma^z_{j,1}\sigma^z_{j+1,1}
    +\sigma^z_{j,2}\sigma^z_{j+1,2}
    +\sigma^z_{j,1}\sigma^z_{j,2})$
   \newline $\displaystyle{}-3\sum_{j,a=1,2}\sigma^x_{j,a}$
 & $L\ge4$ rungs & $0.6368$ & $2.662$\\
Mixed-field Ising
 & $-\sum_j\sigma_j^z\sigma_{j+1}^z
    -\tfrac14\sum_j\sigma_j^x-\tfrac1{20}\sum_j\sigma_j^z$
 & $L\ge6$ sites & $0.3922$ & $1.190$\\
Heisenberg, $S=1$
 & $\sum_j\bm S_j\cdot\bm S_{j+1}$
 & $L\ge36$ sites & $0.0729$ & $0.410$\\
\bottomrule
\end{tabular}
\end{table*}

\emph{Formal verification.---}
The marker $\leanmark{https://sirui-lu.com/CertifiedGaps/docs/CertifiedGaps/Ring/Gap.html}$
denotes a formalized implication with its stated assumptions.
The analytic criterion and recurrence are formalized in Lean.
The path-space kernel construction and trace identities needed for
the original proof are not yet available in mathlib. We therefore
use finite matrices. A positive transfer matrix at finite
$M$ gives closed scalar inequalities; a controlled product limit
at fixed $L,\beta$ passes them to exact quantities before
iteration%
~\leanmark{https://sirui-lu.com/CertifiedGaps/docs/CertifiedGaps/Ring/Gap.html\#CertifiedGaps.RPRing.spatialMomentBounds}.

We prove the gap theorem in Lean for general reflection-positive
Hamiltonians and apply the same theorem to the models in
Table~\ref{tab:models}. We provide the proofs as a
\href{https://github.com/LionSR/CertifiedGaps}{library}
so that others can reproduce the results and add models.

To obtain a complete Lean proof, we must also verify the finite
calculations in Lean. We have done so for the mixed-field Ising model
at $(g,h)=(1/4,1/20)$. For spin one, we check the MPS contractions
in Lean and supply the remaining externally certified finite bounds
as hypotheses.

\emph{Discussion.---}
Using diagrams and visual summaries~\cite{Tufte1990Envisioning} to expose
the structure of the proof let us certify gaps for seven other spin-chain
and ladder models (Table~\ref{tab:models}), and improve the spin-1 bound
using the existing finite inputs. We also explain the boundary-field
proof for open chains using diagrams~\cite{OpenAI2026Boundaryfield,SM}.

A next target is the spin-2 Heisenberg chain, whose smaller gap and
longer correlation length make the finite inputs harder to
certify~\cite{Todo2001Cluster}. In two spatial dimensions, an analogous
argument would track three purities and require additional spatial
inequalities as well as certified finite inputs. Establishing these
inputs for the honeycomb AKLT model~\cite{Lemm2020Existence,Pomata2020Demonstrating}
remains open for this method.

The OpenAI proof combines simple ideas with a substantial computation.
This work also explores numerical certification in Lean and calls for
improvements in the numerical efficiency of proof assistants, making them comparable with
symbolic computation systems such as Mathematica and numerical linear algebra engines such as those in MATLAB and Julia.
The fully checked Ising seed shows that the
finite calculation can be included in the formal proof; the larger
spin-1 calculations provide a further test of how much computation
can be checked directly in the kernel.
